% Generated by scripts/assemble.py; corrections belong in source/corrections.json.
\examhead{174}{西南财经大学 2025 年数学分析试卷}{164}
% p165-b20; source PDF page 165
\blockanchor{p165-b20}
\question{1}求极限 $\lim_{x\to +\infty}e^{-x}\left(1 + \frac{1}{x}\right)^{x^2}$.\par
% p165-b21; source PDF page 165
\blockanchor{p165-b21}
\question{2}判断下面命题是否正确，并证明你的结论.\par
% p165-b22; source PDF page 165
\blockanchor{p165-b22}
函数 $f(x)$ 在 $(a,b)$ 内处处有导数 $f'(x)$，则 $(a,b)$ 中的点不可能为 $f'(x)$ 的第一类间断点.\par
% p165-b23; source PDF page 165
\par\Needspace{6\baselineskip}
\blockanchor{p165-b23}
\question{3}函数 $f(x)$ 在 $[a, b] (ab > 0)$ 上连续，在 $(a, b)$ 上可微. 证明：存在 $c \in (a, b)$，使得\par
\nopagebreak[4]
% p165-b24; source PDF page 165
\blockanchor{p165-b24}
\displayformula{\frac {1}{b - a} \left| \begin{array}{c c} b & a \\ f (b) & f (a) \end{array} \right| = f (c) - c f ^ {\prime} (c).}
% p166-b00; source PDF page 166
\blockanchor{p166-b00}
\question{4}判断积分 $\int_{1}^{+\infty}\frac{\sin^2x}{x^p(x^p + \sin x)}\mathrm{d}x (p > 0)$ 的敛散性.\par
% p166-b01; source PDF page 166
\blockanchor{p166-b01}
\question{5}判断级数 $\sum_{n=2}^{\infty} \frac{n \cos(n\pi)}{\sqrt{n^3 - 2n + 1}}$ 是否收敛, 并指出是绝对收敛还是条件收敛.\par
% p166-b02; source PDF page 166
\par\Needspace{6\baselineskip}
\blockanchor{p166-b02}
\question{6}已知 $f(x, y) = \frac{1}{xy}$, $r = \sqrt{x^2 + y^2}$, $k > 1$.\par
\nopagebreak[4]
% p166-b03; source PDF page 166
\blockanchor{p166-b03}
\displayformula{D _ {1} = \left\{(x, y) \mid \frac {1}{k} x <   y <   k x \right\}, D _ {2} = \{(x, y) \mid x > 0, y > 0 \}.}
% p166-b04; source PDF page 166
\blockanchor{p166-b04}
当 $i = 1,2$ 时，分别判断 $\lim_{\substack{r\to +\infty \\ (x,y)\in D_i}}f(x,y)$ 是否存在.\par
% p166-b05; source PDF page 166
\par\Needspace{6\baselineskip}
\blockanchor{p166-b05}
\question{7}计算二重积分\par
\nopagebreak[4]
% p166-b06; source PDF page 166
\blockanchor{p166-b06}
\displayformula{\iint_ {x ^ {2} + y ^ {2} \leq \frac {3}{16}} \min \left\{\sqrt {\frac {3}{16} - x ^ {2} - y ^ {2}}, 2 (x ^ {2} + y ^ {2}) \right\} d x d y.}
% p166-b07; source PDF page 166
\par\Needspace{6\baselineskip}
\blockanchor{p166-b07}
\question{8}已知\par
\nopagebreak[4]
% p166-b08; source PDF page 166
\blockanchor{p166-b08}
\displayformula{f (x, y) = \left\{ \begin{array}{l l} \frac {x ^ {m} y ^ {m}}{(x ^ {2} + y ^ {2}) ^ {n}} \arcsin (x ^ {2} + y ^ {2}), & (x, y) \neq (0, 0); \\ 0, & (x, y) = (0, 0). \end{array} \right.}
% p166-b09; source PDF page 166
\blockanchor{p166-b09}
判断 $f(x,y)$ 在 $(0,0)$ 处的偏导数是否存在，是否可微.\par
